In this section, we will examine the outcomes of the GAN-generated methods, with a particular focus on evaluating the image quality, diversity, and effectiveness of the generated images on edge cases. The primary objective of this analysis is to highlight the common issues that can arise when using GANs as an augmentation technique and to compare the performance of the three architectures with one another.

Because of their complex architecture and the extensive number of parameters that must be adjusted, Generative Adversarial Networks (GANs) often necessitate a significant amount of time to train. Additionally, the training process can be computationally intensive, requiring high-end hardware and substantial computational resources. Table \ref{tab:gan_comparison} provides an overview of the number of trainable parameters, training time, and epochs required for each GAN employed in this thesis.

\begin{table}[h]
\centering
\begin{tabular}{|c|c|c|c|}
\hline
\textbf{Model} & \textbf{Number of trainable parameters} & \textbf{Epochs} & \textbf{Training Time} \\
\hline
Vanilla GAN & 405574756 & 10000 & 6h 58m 24s \\
WGAN & 743760999 & 5000  & 1d 5h 31m \\
WGAN-GP & 7645523 & 3000 & 10h 4m 50s \\
\hline
\end{tabular}
\caption{Comparison of GAN models}
\label{tab:gan_comparison}
\end{table}

\subsubsection{Image quality} 
\begin{figure}[!h]
	 \centerline{\includegraphics[width=1\textwidth]{src/Results/gan/samples_gan.pdf}}
	 {\caption{Samples of generated images from the conditional GANs}\label{fig:gan_samples}}
\end{figure}
The samples of generated images from three different types of GANs are displayed in Figure \ref{fig:gan_samples}. Each row represents an example of an edge class, such as Schnatterente, Graugans, Wasserlinie mit Großalgen, and Schmarotzerraubmöwe. The first four columns show the input images used to train the different GAN models. 

The images generated from Vanilla GAN appear to be the most realistic compared to other types of GANs. However, it is also observed that Vanilla GAN is prone to getting trapped in a local optima, resulting in generated images that closely resemble a specific real image, as seen in the case of the Schnatterente class. 

As illustrated with small seeming random patches in WGAN, when trained on a small dataset, WGAN did not have enough information to learn the underlying distribution of the data and therefore generate unrecognizable images. This is in contrast to vanilla GAN, which may perform better on smaller datasets but may suffer from mode collapse. 

Wasserstein GAN with gradient penalty has been demonstrated to produce generated samples that are more similar to the distribution of the training data in comparison to Vanilla WGAN. Although the generated images still exhibit a deficiency in high-resolution details, they share some similarities with the objects in the dataset. For instance, the class of Schnatterente features white dots and yellow patches, while the class of Wasserlinie mit Großalgen displays a similar shade of blue as shown in \ref{fig:gan_samples}. This observation highlights the issue of hypersensitivity to hyperparameters in WGAN. In the absence of hyperparameter tuning for the model, it failed to converge rapidly.

\subsubsection{Diversity}
Figure \ref{fig:gan_divers} displays different samples generated by 3 types of GANs with varying noise inputs across 3 classes. It is evident that Vanilla GAN suffers from the mode collapse problem, as it only produces one repeated image with slight variations in brightness and hue. On the other hand, WGAN and WGAN with GP successfully generate diverse samples.

\begin{figure}[!h]
	 \centerline{\includegraphics[width=1\textwidth]{src/Results/gan/divers_gan.pdf}}
	 {\caption{Variety of generated pictures in three GANs}\label{fig:gan_divers}}
\end{figure}

\subsubsection{Performance} 

\begin{figure}[!h]
\hfill
\subfigure[F1-Score Top-3 for all classes \label{fig:overall_gan}]{\includegraphics[width=7cm]{src/Results/gan/overall_gan.png}}
\hfill
\subfigure[F1-Score Top-3 for only edge classes \label{fig:edge_gan}]{\includegraphics[width=7cm]{src/Results/gan/edge_gan.png}}
\hfill
\caption{F1-Score (a) for all classes and (b) for edge classes}
\label{fig:gan_performance}
\end{figure}

% \begin{figure}[!h]
% 	 \centerline{\includegraphics[width=0.7\textwidth]}
% 	 {\caption{}}
% \end{figure}
Figure \ref{fig:overall_gan} depicts the F1-Score of EfficientNet when trained with additional GAN-generated data, while Figure \ref{fig:edge_gan} displays the F1-Score for edge cases. Overall, the overall classification performance decreases with the use of GAN-generated data. One possible explanation is that the GAN-generated images may contain noise or inconsistencies that are absent in the original dataset, thereby confusing the model and making it more difficult to generalize to new examples. Additionally, GAN-generated images may not cover the entire distribution of the original dataset, leading to overfitting and poor generalization performance.



WGAN with GP causes the largest decrease in F1-Score Top-3, with a decrease of nearly 5\%. On the other hand, using data from Vanilla GAN leads to the smallest decrease in performance, with a decrease of only 1.8\%. For the edge classes, both WGAN and Vanilla GAN show an improvement in F1-Score Top-3. WGAN results in a larger improvement of 0.013, while Vanilla GAN yields a smaller improvement of 0.0014. However, it is important to consider the overall decrease in performance when using GAN-generated data. The primary reason for the improvement in edge cases may be the introduction of inconsistencies and new features, which the network then learns to recognize. However, this improvement may not outweigh the decrease in performance seen in the rest of the dataset. In the case of WGAN with GP, the overall decrease in performance is 0.0355, which may indicate that the network is overfitting to the GAN-generated data. In the case of Vanilla GAN, the overall decrease in performance is 0.0205, which is still significant but less than that of WGAN with GP.




